\documentclass[12pt] {amsart}
\newcommand{\OO}{{\mathcal {O}}}
\begin{document}
\title{Research of Hacon Christopher}
\maketitle

My research is mostly focoused on the birational geometry of algebraic varieties. I have been particularly interested in the birational geometry of irregular varieties and of linear systems on abelian varieties, in the birational geometry of surfaces and 3-folds and recently in natural questions on the birational geometry of higher dimensional varieties which naturally arise in teh context of the minimal model program. (The references below are from the attached list of publications. Many of the cited works are joint publications. The co-authors are listed in the attached list of publications.)

\bigskip
\noindent{\bf Singularities of divisors on abelian varieties:} In [3], [4] and [23], we complete and extend the results of Koll\'ar and of Ein-Lazarsfeld on singularities of Theta divisors, to the singularities of divisors in $|mL|$ where $L$ is an ample divisor on an abelian variety with $h^0(\OO (L))\leq 3$.
 

\bigskip
\noindent{\bf Birational geometry of irregular varieties:}
 
In [5] and [7] we prove that (as conjectured by Koll\'ar), any variety with
$P_2(X)=1$ and $q(X)=\dim X$ is birational to an abelian variety.
In [11], [14] and [18] we show that one can classify all varieties with
$P_3(X)\leq 4$ and $q(X)=\dim X$.
In [7], we give some natural results closely connected to a conjecture of Ueno (Conjecture K). Further evidence towards this conjecture is given in [17].
In [11], we also give a birational characterization of Theta divisors. This
characterization is surprisingly sharp In fact it leads to a solution of a 
long standing problem of classifying surfaces with $p_g=q=3$ (this is done in in [10]).
In [19], we construct a counter-example to the corresponding natural conjecture for surfaces with $p_g=q=2$.
In [20], we show that one can give a birational characterization of
products of curves of genus $2$ thus extending a result on surfaces by Beauville to the case of arbitrary dimension.
In [22], we give some partial answers to a natural conjecture which states that on a variety of general type, any holomorphic $1$-form has a zero. (We adress the case in which the abelian variety is simple and the case in which the ambient variety is smooth and minimal.)
In [17] we give a proof of a natural conjecture of Green and Lazarsfeld. This extends and gives a better understanding of their results on generic vanishing theorems. A similar result was obtained independantely and by completely different methods by G. Pareschi.
In [9] we give a unified treatment of the results of the generic vanishing reaults of Simpson, Arapura, Green-Lazarsfels and others. We also give applications of these results to higher direct images of dualizing sheaves and to the divisibility of certain divisors on an abelian variety.
In [15], we recover an analogue (but stronger statement) to the Fujita conjecture for varieties of maximal Albanese dimension. This leads to almost optimal results on the pluricanonical maps of varieties of general type and maximal Albanese dimension. G. Pareschi and M. Popa, have also independently obtained similar results via similar methods.

\bigskip
\noindent{\bf Birational geometry of 3-folds:}
Let $X$ be a Gorenstein minimal complex projective 3-fold of general type, then the geometry of $X$ behaves quite similarly to the birational geometry of surfaces of general type. In particular, the geometry of the pluricanonical maps
is well understood. For example, a recent result of Chen, Chen and Zhang, states that the $5$-th canonical map is birational. In [21], we show that if $X$
is a $3$-fold as above then $\chi (\omega _X)\leq 2p_g$. It follows that
the degree of the canonical map is bounded. Meng Chen has shown that if
the canonical map has fibers of dimension $1,2$, then their genus $g,p_g$ is
bounded from above. In [27], we further refine these results which are an analogue of the results of Beauville for the canonical maps of surfaces of general type.

In a different direction, I expect that one should be able to prove analogous results for all but finitely many families of $3$-folds, or equivalently buy a resut of Tsuji (see [24]) for all $3$-folds with sufficiently big canonical volume. For example, I expect that in the near future, my student G. todorov, will show that: ``for alll but finitely many families of $3$-folds of general type, one has $P_2>0$ and $|5K_X|$ gives a birational map.''

\bigskip
\noindent{\bf Birational geometry of higher dimensional varieties:}

In [24] we prove a new extension theorem that shows that certain pluricanonical forms extend from a divisor to the ambient variety. This result is based on techniques of Siu and Kawamata and is inspired by an observation of Tsuji.
It differs from Tsuji's statement (also recovered by Takayama), in that it readily applies to the context of the log minimal model program.
In [24], we apply this result to recover the following result of Tsuji:
``For every positive integer $n>0$, there is a fixed constant $r=r_n$, such that if $X$ is a variety of general type and dimension $n$, then the $r$-th pluricanonical map is birational.''
In [25] we recover a conjecture of Shokurov which in particular implies that the fibers of a resolution of a variety with divisorial log terminal singularities are rationally chain connected.
In [26], we prove the existence of flips in dimension n, contingent on the termination of real flips in dimension n-1. In particular, since Shokurov (building on work of Mori, Kawamata, Koll\'ar and others) has shown the existence and termination of $3$-fold real flips, this gives a new and substantially simplified proof of the existence of $4$-fold flips.

\bigskip
\noindent{\bf Other research topics:}
In [1] and [2], we study the seshadri constants of vector bundles. We give a counterexample to a conjecture of Ballico and of Beltrametti-Sommese-Schneider which roughly speaking states that the positivity properties of an ample vector bundle should be closely related to the positivity properties of very ample vector bundles. We also give some affirmative answers to several related questions that extend the known results for line bundles.

In [8], [16], we generalize the description of the subring of the coordinate ring for the $n \times r$ matrices $k [ M_{nr}]$, $r < n$, consisting of all the invariant elements with respect to the right action of the special linear group ${\rm SL}_r(k)$, where $k$ is an algebraically closed field of characteristic $0$ to the case where $k [ M_{nr}]$ is replaced by the quantum matrix bialgebra and ${\rm SL}_r(k)$ by the quantum special linear group.


\end{document} 

